\documentclass[10pt,a4paper]{amsart}
\usepackage[margin=19mm]{geometry}
\usepackage{amsmath,amssymb,mathtools}
\usepackage[scr=rsfs,cal=euler]{mathalfa}
\usepackage{xfrac}
\usepackage[unicode,hidelinks,bookmarks=false]{hyperref}
\newcommand{\ie}{i.e.\ }
\newcommand{\cf}{cf.\ }
\newcommand{\resp}{respectively\ }
\newcommand{\Subsection}{\subsection}
\providecommand{\nobreakdash}{\nobreak\hbox{-}}
\setlength{\emergencystretch}{2em}
\multlinegap0pt
\usepackage{amssymb}
\usepackage{enumerate}
\usepackage{tikz}
\usepackage{mathtools}
\usepackage{dsfont}
\usepackage{csquotes}
\newtheorem{corollary}{Corollary}
\newcommand{\Dom}{\mathrm{Dom}}
\title{Obata's rigidity theorem in free probability}
\author{Charles-Philippe Diez}
\date{Source: arXiv:2603.05466v1 (CC BY 4.0)}
\begin{document}
\maketitle

\noindent\emph{Source and licence.} Obata's rigidity theorem in free probability. Version arXiv:2603.05466v1 (CC BY 4.0), distributed by arXiv under the Creative Commons Attribution 4.0 International licence, http://creativecommons.org/licenses/by/4.0/, as recorded in the arXiv licence metadata for this exact version and shown on the abstract page https://arxiv.org/abs/2603.05466v1. Full licence notice: \texttt{licenses/arxiv-2603.05466-CC-BY-4.0.txt}.\par\smallskip

\noindent\textit{Editorial scope.} Counted unit: the article's corollary cor1 (source0.tex lines 2560--2562). The article's complete original proof of it is delivered with it (source0.tex lines 2564--2629, one \texttt{proof} environment of 66 lines). No mathematics is written, repaired or re-derived: the statement and the proof are the article's own lines, in the article's own notation, and no retained line is substituted or re-flowed.  No further source-local context is needed: every cross-reference of every retained line resolves inside the two delivered spans.  Nothing was omitted from any retained span, so the package carries no omission record.  No retained line contains a citation, so the wrapper carries no bibliography and no bibliography entry.  No retained line contains a figure environment, an image-inclusion command or drawing code, so no image is delivered and the package declares no asset dependency.  Editorial formatting only, in addition: an AMS \texttt{amsart} preamble that restates the article's own declarations and macros used by the retained lines, namely the environments \texttt{corollary} on separate counters, together with the standard packages those lines need.  The retained environments print in this document as Corollary 1 in order of appearance, whereas the article prints the same items with the numbering of its own full document.  The complete native source of the article as distributed in the arXiv e-print for this exact version is bundled unchanged as \texttt{sources/195804/source0.tex}.
\begin{corollary}\label{cor1}
Under the curvature assumption $CD(1,\infty)$, if $f\in \Dom(\Delta)\cap L^2_0(M)$ satisfies $\Delta f=f$ and $\|f\|_{L^2(M)}^2=\mathcal{E}(f)$, then so do $\Re f$ and $\Im f$. In particular, without loss of generality one may assume $f=f^\ast:=Jf$.
\end{corollary}

\begin{proof}
Write $f=u+iv$ with $u=\Re f$, $v=\Im f$. Since $\tau(f)=0$, also $\tau(u)=\tau(v)=0$.  
We also have the $\ast$–stability of $\Dom(\Delta)$ and complex linearity, checked using the \emph{real} property of free difference quotients. Indeed, fix $x\in \Dom(\Delta)$, $y\in \Dom( \partial)$, then
\begin{eqnarray}
    \langle \Delta (x^*),y\rangle_{L^2(M)}&=&\sum_{i=1}^n \langle \bar \partial_i (x^*),\bar\partial_i y\rangle_{L^2(M)\bar\otimes L^2(M^{op})}\nonumber\\
    &=&\sum_{i=1}^n \langle (\bar \partial_i x)^{\dagger},\bar\partial_i y\rangle_{L^2(M)\bar\otimes L^2(M^{op})}\nonumber\\
    &=& \sum_{i=1}^n \langle (\bar \partial_i y)^{\dagger},\bar\partial_i x\rangle_{L^2(M)\bar\otimes L^2(M^{op})}\nonumber\\
    &=& \sum_{i=1}^n \langle \bar \partial_i (y^*),\bar\partial_i x\rangle_{L^2(M)\bar\otimes L^2(M^{op})}\nonumber\\
    &=& \langle y^*,\Delta x\rangle_{L^2(M)}=\langle (\Delta x)^*,y\rangle_{L^2(M)}
\end{eqnarray}
where we used for $a,b\in L^2(M)\bar\otimes L^2(M^{op})$ that $\langle a,b\rangle_{L^2(M)\bar\otimes L^2(M^{op})}=\langle b^{\dagger},a^{\dagger}\rangle_{L^2(M)\bar\otimes L^2(M^{op})}$. Hence $x^*\in \Dom(\Delta)$, and $\Delta(x^*)=(\Delta(x))^*:=J\Delta(x)$.
\bigbreak

Hence, we have $\Delta(f^*)=f^*$, and we get,

\[
\Delta u=\Re(\Delta f)=\Re(f)=u,\qquad
\Delta v=\Im(\Delta f)=\Im(f)=v.
\]


For the norms, traciality of $\tau$ gives


\[
\|f\|_{L^2(M)}^2=\|u\|_{L^2(M)}^2+\|v\|_{L^2(M)}^2.
\]


For the Dirichlet energies, the \emph{real} property for the free difference quotients and orthogonality yield


\[
\mathcal{E}(f)=\mathcal{E}(u)+\mathcal{E}(v).
\]


By hypothesis $\|f\|^2_{L^2(M)}=\mathcal{E}(f)$, hence


\[
\|u\|^2_{L^2(M)}+\|v\|^2_{L^2(M)}=\mathcal{E}(u)+\mathcal{E}(v).
\]


Therefore,


\[
(\|u\|^2_{L^2(M)}-\mathcal{E}(u))+(\|v\|^2_{L^2(M)}-\mathcal{E}(v))=0,
\]


By the Poincaré inequality under $CD(1,\infty)$, we have
$\|u\|_{L^2(M)}^2\le \mathcal{E}(u)$ and $\|v\|_{L^2(M)}^2\le \mathcal{E}(v)$, so each
difference $\|u\|_2^2-\mathcal{E}(u)$ and $\|v\|_2^2-\mathcal{E}(v)$ is
non-positive. Since their sum is zero, they must both vanish, and hence, forces equality termwise:


\[
\|u\|^2_{L^2(M)}=\mathcal{E}(u),\qquad \|v\|^2_{L^2(M)}=\mathcal{E}(v).
\]

\bigbreak
Thus $u,v$ are centered (eigen)functions saturating the free Poincaré inequality. Hence, we may assume without loss of generality $f=f^\ast$.
\end{proof}

\end{document}
